\documentclass[12pt,twoside,a4paper,chapter=TITLE]{abntex2}

\usepackage[T1]{fontenc}
\usepackage[utf8]{inputenc}
\usepackage[brazil]{babel}
\usepackage{booktabs}
\usepackage[alf]{abntex2cite}

\titulo{Toler\^ancia a Parti\c{c}\~oes de Rede em Protocolos de Consenso Aproximado}
\autor{Jo\~ao da Silva}
\local{S\~ao Paulo}
\data{2026}
\instituicao{Universidade de Ci\^encias e Tecnologia\\Faculdade de Computa\c{c}\~ao\\Programa de P\'os-Gradua\c{c}\~ao em Ci\^encia da Computa\c{c}\~ao}
\tipotrabalho{Disserta\c{c}\~ao}
\preambulo{Disserta\c{c}\~ao apresentada ao Programa de P\'os-Gradua\c{c}\~ao em Ci\^encia da Computa\c{c}\~ao como requisito parcial para a obten\c{c}\~ao do t\'itulo de Mestre.}
\orientador{Prof. Dr. Maria Souza}

\begin{document}

\imprimircapa
\imprimirfolhaderosto*

\begin{resumo}
Esta disserta\c{c}\~ao estuda a degrada\c{c}\~ao de protocolos de consenso
aproximado sob parti\c{c}\~ao parcial de rede, e prop\~oe um protocolo que
troca um pequeno aumento de lat\^encia por maior vaz\~ao.

 \vspace{\onelineskip}

 \noindent
 \textbf{Palavras-chave}: sistemas distribu\'idos. consenso. parti\c{c}\~ao de rede.
\end{resumo}

\begin{resumo}[\begin{otherlanguage*}{english}Abstract\end{otherlanguage*}]
 \begin{otherlanguage*}{english}
This thesis studies how approximate consensus protocols degrade under
partial network partition, and proposes a protocol that trades a small
latency overhead for improved throughput.

 \vspace{\onelineskip}

 \noindent
 \textbf{Keywords}: distributed systems. consensus. network partition.
 \end{otherlanguage*}
\end{resumo}

\pretextual

\tableofcontents

\textual

\chapter{Introdu\c{c}\~ao}
\label{ch:introducao}

\section{Motiva\c{c}\~ao}

Sistemas distribu\'idos de grande escala t\^em se apoiado cada vez mais
em protocolos de consenso aproximado para equilibrar lat\^encia e
consist\^encia \cite{lamport1998part}. Este trabalho estuda como tais
protocolos se degradam sob parti\c{c}\~ao de rede.

\section{Contribui\c{c}\~oes}

As principais contribui\c{c}\~oes s\~ao:

\begin{enumerate}
  \item Um modelo formal de toler\^ancia a parti\c{c}\~oes parciais.
  \item Uma avalia\c{c}\~ao emp\'irica resumida na Tabela~\ref{tab:resultados}.
  \item Uma implementa\c{c}\~ao de refer\^encia usada para gerar a
    Figura~\ref{fig:latencia}.
\end{enumerate}

\section{Trabalhos Relacionados}

Protocolos de consenso cl\'assicos como Paxos \cite{lamport1998part} e
Raft \cite{ongaro2014search} assumem uma rede totalmente conectada.
Esta disserta\c{c}\~ao relaxa essa hip\'otese.

\chapter{Avalia\c{c}\~ao}
\label{ch:avaliacao}

\section{Configura\c{c}\~ao Experimental}

Os experimentos foram executados em um cluster de cinco n\'os. A
Tabela~\ref{tab:resultados} resume a vaz\~ao sob taxas crescentes de
parti\c{c}\~ao.

\begin{table}[htbp]
  \centering
  \caption{Vaz\~ao (requisi\c{c}\~oes/s) sob taxas simuladas de parti\c{c}\~ao.}
  \label{tab:resultados}
  \begin{tabular}{@{}lrr@{}}
    \toprule
    Taxa de parti\c{c}\~ao & Linha de base & Proposto \\
    \midrule
    0\%  & 12.400 & 12.100 \\
    5\%  & 9.800  & 11.600 \\
    10\% & 6.200  & 10.900 \\
    \bottomrule
  \end{tabular}
  \fonte{Elaborado pelo autor.}
\end{table}

\section{Lat\^encia sob Carga}

A Figura~\ref{fig:latencia} mostra a lat\^encia mediana das
requisi\c{c}\~oes conforme a taxa de parti\c{c}\~ao aumenta.

\begin{figure}[htbp]
  \centering
  \rule{0.6\textwidth}{4cm}
  \caption{Lat\^encia mediana em fun\c{c}\~ao da taxa de parti\c{c}\~ao simulada.}
  \label{fig:latencia}
  \fonte{Elaborado pelo autor.}
\end{figure}

\section{Discuss\~ao}

O protocolo proposto mant\'em uma vaz\~ao maior que a linha de base
assim que a taxa de parti\c{c}\~ao ultrapassa 5\%, ao custo de um
pequeno aumento constante de lat\^encia \cite{ongaro2014search}.


\postextual

\bibliography{referencias}

\end{document}
